\section{The point row and three distinct frames}\label{sec:point-row}

We fix the pairing convention before discussing wall crossing.  Let
\(\mathcal S(Y)\) denote a space of flat sections of the quantum connection on
a domain where the Gamma framing is defined.  Iritani's section associated
with \(E\in K^0_{\mathrm{top}}(Y)\) has the form
\begin{equation}\label{eq:gamma-framed-section}
 s_Y(E)=L_Y(\tau,z)z^{-\mu}z^{c_1(Y)}
 (2\pi)^{-\dim Y/2}\Ghat_Y(2\pi i)^{\deg/2}\operatorname{ch}(E),
\end{equation}
up to the standard choice of normalization.  The flat pairing is
\begin{equation}\label{eq:flat-euler-pairing}
 [s_1,s_2):=(s_1(\tau,e^{-\pi i}z),s_2(\tau,z))_Y,
 \qquad [s_Y(E),s_Y(F))=\chi(E,F).
\end{equation}
See \cite[Section~1]{IritaniGamma}.

\begin{definition}[Gamma point row]\label{def:point-row}
For \(p\in Y\), the Gamma point row is
\begin{equation}\label{eq:gamma-point-row}
 \rrow_{Y,p}(v):=[v,s_Y(\OO_p)).
\end{equation}
On Gamma-framed \(K\)-classes it is ordinary rank:
\begin{equation}\label{eq:point-row-rank}
 \rrow_{Y,p}(s_Y(E))=\chi(E,\OO_p)=\rk(E).
\end{equation}
\end{definition}

Putting the point in the first entry instead would introduce the usual
dimension sign.  We use \eqref{eq:gamma-point-row} throughout.

Three levels of structure occur in the arguments below.
\begin{enumerate}[label=\textup{(\roman*)}]
\item A \emph{formal packet} is a generalized formal-monodromy or
  exponential block of the \(z=0\) formal connection.
\item A \emph{sectorial lift} is a choice of actual flat subspace with that
  formal asymptotic expansion on an admissible sector.
\item The \emph{intrinsic point section} is the Gamma-normalized flat
  section fixed by the large-radius fundamental solution.
\end{enumerate}

A positive-\(z\) gauge can identify formal quantum D-modules without
identifying (iii) across two incompatible Novikov completions.  Likewise,
two sectorial lifts of the same formal packet can differ by a Stokes shear.
The zero or nonzero restriction of \eqref{eq:gamma-point-row} is therefore
not determined by the
formal packet alone.

We write \(\cP_6(Y)\) for the generalized formal-monodromy block with
eigenvalue \(\zetaSix=e^{2\pi i/6}\).  When \(\cP_6(Y)\) appears inside a
point-row Boolean, it denotes a specified sectorial lift of that block; a
different lift is different data.  An empty block has Boolean value zero.

\begin{definition}[Point-row Boolean]\label{def:boolean}
For a sectorially realized formal packet \(\cP\subset\mathcal S(Y)\), put
\begin{equation}\label{eq:point-row-boolean}
 b(Y,\cP)=
 \begin{cases}
 1,&\rrow_{Y,p}|_{\cP}\ne0,\\
 0,&\rrow_{Y,p}|_{\cP}=0.
 \end{cases}
\end{equation}
\end{definition}

The paper's exact wall identities either identify the intrinsic point
section directly or remain within one specified receiver.  Whenever an
additional sectorial realization is required, it is stated as a hypothesis.
